%%%%%%%%%%  scan idx 191  |  folio 178  %%%%%%%%%%
équivalences de catégories :
\[
\begin{gathered}
\mathrm D(S)\longrightarrow\mathrm D(A(S))\\
\mathrm D^{-}(S)_{\coh}\longrightarrow\mathrm D^{-}(A(S))_{\coh}\\
\mathrm D(S)_{\parf}\longrightarrow\mathrm D(A(S))_{\parf},
\end{gathered}
\]
les notions de pseudo-cohérence (resp. perfection) dans $\mathrm D(A(S))$ étant entendues relativement à la catégories des modules libres (resp. projectifs) de type fini. Il définit aussi une équivalence de la sous-catégorie pleine de $\mathrm D(S)$ formée des objets $n$-pseudo-cohérents (resp. d'amplitude parfaite contenue dans $[m,n]$) sur la sous-catégorie pleine analogue de $\mathrm D(A(S))$.
\end{enumerate}

\emph{Preuve.} La première assertion de a) résulte de (2.1.2 a)) appliqué en prenant pour $S$ le topos des faisceaux d'ensembles sur $S$, et $\mathcal C'_S=\Mod(S)$ (compte tenu de (2.3.1)). La seconde découle de (1.1 a)) (resp. (1.2 a))) avec $\mathcal C=\Mod$ (catégorie fibrée des $A$-Modules sur des ouverts variables), $\mathcal C_0=\Lib$ (resp. $\Loclib$), la condition (ii) de (1.1) étant vérifiée d'après (2.3.1) et (2.1.1.b)). Pour l'assertion b), on applique (1.1 b),c)) (avec $\mathcal C$ et $\mathcal C_0$ comme précédemment). L'assertion c) résulte de (1.2.1 b),c),d)), où $\mathcal C=\Mod$ et $\mathcal C_0=\Lib$. La première assertion de d) découle, grâce à (2.3.1), de (2.1.2 b) (ii)) où l'on fait $\mathcal C'_S=\mathcal C_S=\Mod(S)$. Les autres assertions de d) se déduisent de la précédente, compte tenu de b) et c), et du fait que le foncteur $\Gamma_S$ transforme Modules libres (resp.) facteurs directs de libres) de type fini en modules libres (resp. projectifs de type fini).

\subsection*{2.4. Espaces de Stein}

\textbf{Définition 2.4.1.} Soit $S$ un topos localement annelé (I 2.15.1 (ii)), d'objet final $S$. On dit que $S$ est de Stein si le faisceau structural $\mathcal O_S$ est cohérent (I 3.1) et si l'on a les propriétés suivantes :

«~\emph{théorème} A~» : tout $\mathcal O_S$-Module cohérent (I 3.1) est engendré par ses sections,

«~\emph{théorème} B~» : pour tout $\mathcal O_S$-Module cohérent $F$, on a $H^i(S,F)=0$ pour $i>0$.

